\documentclass[10pt,a4paper]{article}
\usepackage[utf8]{inputenc}
\usepackage[T1]{fontenc}
\usepackage[brazilian]{babel}
\usepackage[left=0.55in,top=0.45in,right=0.55in,bottom=0.45in]{geometry}
\usepackage{titlesec}
\usepackage{enumitem}
\usepackage{hyperref}
\usepackage{xcolor}
\usepackage{charter}

% Section formatting
\titleformat{\section}{\large\bfseries\uppercase}{}{0em}{}[\titlerule]
\titlespacing{\section}{0pt}{8pt}{4pt}

% Itemize formatting
\setlist[itemize]{leftmargin=0.15in, labelsep=0.08in, topsep=1pt, itemsep=2pt, parsep=0pt}

\pagestyle{empty}

\begin{document}

% Cabeçalho
\begin{center}
    {\huge \textbf{João Victor Mendes}} \\
    \vspace{3pt}
    {\large Production Support Engineer (PSE) / SRE | Cloud · Observabilidade · Database Reliability} \\
    \vspace{2pt}
    \href{https://linkedin.com/in/mendes-victor}{linkedin.com/in/mendes-victor} | \href{https://github.com/victormends}{github.com/victormends} | \href{mailto:j.victor8@hotmail.com}{j.victor8@hotmail.com} | +55 (41) 99593-8387 \\
    União da Vitória, PR, Brasil (UTC-3)
\end{center}

% Resumo Profissional
\section{Resumo Profissional}
\textbf{Medalha de Prata, Olimpíada Brasileira de Matemática (OBMEP, Top 0,01\% entre 18M+ participantes).} Traduz rigor analítico extremo em resposta estrutural a incidentes, observabilidade e engenharia de confiabilidade de plataformas de alta escala (IA e Fintech). Especialista em diagnóstico profundo de sistemas distribuídos, mitigação de gargalos de saturação (AWS Aurora MySQL, Kubernetes, Airflow) e eliminação estrutural de ruído em telemetria, com histórico de redução de mais de 70\% em ingestão e liderança de post-mortems críticos de produção.

% Habilidades Técnicas
\section{Habilidades Técnicas}
\begin{itemize}
    \item \textbf{Cloud \& Orquestração:} AWS (Aurora MySQL/RDS, EKS, ECS, SQS, SSM, CloudWatch, Secrets Manager, SES), Kubernetes (Karpenter, ArgoCD, KEDA, Topology Spread Constraints), Docker.
    \item \textbf{Observabilidade \& SRE:} New Relic (NRQL, NerdGraph API, baseline/alerting), OpenTelemetry (OTel Collector), Resposta a Incidentes P1/P2, RCA / Pós-Mortem (ITIL), Engenharia de Detecção \& Runbooks, SLA.
    \item \textbf{Bancos de Dados \& Storage:} AWS Aurora MySQL (Performance Insights, connection pooling, parameter groups), PostgreSQL (WAL, replication slots, bloat, EXPLAIN ANALYZE, pg\_class), AWS Redshift, Redis / Valkey.
    \item \textbf{IaC, Automação \& Linguagens:} Terraform, Terragrunt, Atlantis, Python, Bash, PowerShell, C\#/.NET, SQL, Git.
    \item \textbf{Idiomas:} Inglês (Proficiência Profissional Plena Certificada B2, BRASAS), Português (Nativo).
\end{itemize}

% Experiência Profissional
\section{Experiência Profissional}

\textbf{Monest} | Fintech de Recuperação de Crédito e Plataforma de IA \hfill Ago 2026 -- Atual \\
\textit{Production Support Engineer (PSE) / SRE} \hfill Remoto, Brasil

Atuação na confiabilidade, observabilidade e infraestrutura de plataforma distribuída de alta escala, sustentando agentes de IA (MIA), discadores VoIP (Jambonz) e orquestração de microsserviços em nuvem.
\begin{itemize}[topsep=3pt]
    \item \textbf{Diagnóstico de Saturação de Banco \& Capacidade de Pool (Aurora MySQL):} Liderou investigação de saturação crítica no Aurora (writer a 98,8\% CPU e teto de 1.000 conexões esgotado com 383 jobs caídos em 27 empresas nas filas do agente de IA); quantificou o fan-out do pool de conexões entre microsserviços (23 instâncias $\times$ 50 = 1.150 conexões potenciais vs. teto rígido de 1.000) e queries sem loteamento, destravando a priorização de RDS Proxy e desacoplamento de parameter groups dedicados.
    \item \textbf{Observabilidade Zero Noise \& Falhas Estruturais de Alerta (New Relic / NerdGraph):} Auditou e classificou 467 conditions da conta via NerdGraph API; isolou política sem faceta gerando \textasciitilde{}4.900 incidentes/mês (\textasciitilde{}40\% do volume total da conta) e descobriu falha matemática em alertas de baseline (\texttt{ALL} com duração superior à janela em crons esparsos), prevenindo notificações silenciosas e formalizando a correção na RFC de governança de observabilidade da companhia.
    \item \textbf{Resposta a Outage de Pipeline \& Engenharia de Detecção (Airflow / Docker / SSM):} Diagnosticou e liderou pós-mortem de incidente de 3h45 no cluster Airflow causado por estouro de disco (ENOSPC em exportação de 108k gravações); contornou bloqueio de SSH via AWS Systems Manager (SSM) sem privilégios de admin, formulou recomendação de alerta deadman no scheduler reduzindo MTTD de 2h02 para \textasciitilde{}10 minutos e corrigiu falha silenciosa em regras de alerta de disco via IaC, onde médias multi-mount mascaravam saturação de 99,89\% em volumes raiz.
    \item \textbf{Resolução de Outage de Agendador \& Persistência (Cronicle / Systemd / IaC):} Identificou a causa raiz de um apagão de 13h no scheduler (154 execuções perdidas em 4 pilares operacionais) não detectado por 472 conditions devido a health checks que validavam o nginx em vez da execução de jobs; projetou a persistência do host via units de serviço systemd, montagem resiliente com \texttt{nofail} no fstab e inicialização via pm2 no boot, validado e aplicado em produção em janela segura de 179 minutos sem indisponibilidade.
    \item \textbf{Resiliência K8s \& Eliminação de SPOF em Stack de IA (EKS / Karpenter / ArgoCD):} Solucionou outage de 21 minutos na stack de observabilidade de IA (Langfuse/Valkey) após interrupção de nó Spot derrubar pods de web, worker e Valkey de réplica única; reestruturou manifests via ArgoCD com \texttt{topologySpreadConstraints} estrito (\texttt{whenUnsatisfiable: DoNotSchedule}), eliminando empacotamento em nó único pelo Karpenter e garantindo distribuição cross-node em produção.
    \item \textbf{Governança de Alertas RDS \& Detecção de Incidentes (Terraform / New Relic):} Codificou 8 políticas de monitoramento do Aurora no Terraform com limiares calibrados sobre 336 buckets de 30 minutos em 7 dias; identificou e corrigiu falha estrutural de deduplicação que mantinha alertas de carga continuamente abertos (em breach durante 37,8\% do tempo no writer e 53,6\% na réplica) e suprimia avisos subsequentes; os novos alertas detectaram anomalias críticas de conexão 4 dias após o rollout.
\end{itemize}

\newpage

\textbf{Pedroso Automação} | ERP SaaS de conformidade fiscal B2B \hfill Junho 2024 -- Ago 2026 \\
\textit{Engenheiro de Suporte Técnico | Recuperação de Infraestrutura L2/L3} \hfill União da Vitória, PR

Suporte L2/L3 direto para 500+ ambientes corporativos (1.500+ usuários finais). Manteve independência total em rotações solo, resolvendo incidentes P1/P2 sem escalada por meio de análise estrutural profunda.
\begin{itemize}[topsep=3pt]
    \item \textbf{Eliminou a classe de risco de retenção de WAL} na frota: construiu um monitor PowerShell agendado via Agendador de Tarefas (PG12) que detecta slots de replicação órfãos inativos e os remove quando a retenção ultrapassa 2GB ou a inatividade supera 30 minutos, com alertas no Visualizador de Eventos. Em instâncias PG13+, complementado com \texttt{max\_slot\_wal\_keep\_size} como limite rígido no servidor.
    \item \textbf{Automatizou a recuperação de cluster com 30 bancos de dados} desenvolvendo script PowerShell que paraleliza verificações de estado, limpa arquivos \texttt{postmaster.pid} obsoletos e monitora recuperação WAL na inicialização; reduziu o tempo de recuperação de horas para menos de 5 minutos.
    \item \textbf{Padronizou a implantação do PostgreSQL no Windows} com script de instalação ponta a ponta, configuração de serviço e hardening de runtime; reduziu o tempo de configuração manual de \textasciitilde{}2 horas para menos de 7 minutos nos ambientes da frota.
    \item \textbf{Diagnosticou bloat de disco de 1GB/mês} por meio de \texttt{pg\_class} e \texttt{pg\_toast}; rastreou a causa raiz a caminhos de arquivos criptografados serializados como milhares de caracteres por entrada de log de auditoria; reduziu o tempo de restore em 60\% (17 min para 6:58) via \texttt{pg\_restore -j 4}.
    \item \textbf{Construiu pipeline ETL de 70M de registros} com redução de 60\% no tempo de processamento sobre dataset governamental de 15GB+; estendido com classificação tributária de CNPJs (MEI, Simples Nacional, Normal, inválido), auditoria de regime fiscal via XML (NF-e/NFC-e/NFS-e/CT-e) e staging em PostgreSQL, expondo o top 0,001\% dos leads comerciais ativos.
    \item \textbf{Protegeu R\$500k em estoque ativo (\textasciitilde{}\$100k USD)} ao fazer engenharia reversa de um banco Firebird legado sem documentação; reprocessou 20 anos de changelogs de transações para reconstruir o estado completo do inventário com zero perda de dados, entregue no mesmo dia em que o esquema foi visto pela primeira vez.
    \item \textbf{Reduziu diagnóstico fiscal de 30+ min para menos de 5 min} desenvolvendo fluxo estruturado de triagem assistido por LLM: extração direta de regras tributárias de rejeição SEFAZ com injeção de contexto de banco de dados e sistema, eliminando toda uma classe de escaladas para desenvolvimento.
    \item \textbf{Estabeleceu runbooks e cultura de mentoria técnica}; produziu 20+ guias em vídeo e pós-mortems ITIL, reduzindo a dependência de disponibilidade sênior e acelerando o onboarding de novos engenheiros.
\end{itemize}

\vspace{6pt}
\textbf{Girafa Comunicação Interativa} | Agência Digital \hfill Set 2022 -- Abr 2024 \\
\textit{Analista de Infraestrutura Cloud \& Operações Web} \hfill Remoto, Brasil

Operou infraestrutura Cloud e operações web para 70+ clientes corporativos (AWS Route 53, Registro.br, Locaweb).
\begin{itemize}[topsep=3pt]
    \item \textbf{Operação e Resolução de Incidentes Cloud:} Gerenciou DNS corporativo, WAF e zonas no AWS Route 53 e CloudFront; atuou como liaison internacional com registradores externos (GoDaddy, AWS, Namecheap) em migrações complexas de infraestrutura e isolamento de falhas de propagação/cache.
    \item \textbf{Investigação de Incidentes de Segurança:} Auditou históricos de Git e logs de acesso para identificar brechas de segurança; revogou sessões comprometidas e realizou rotação de credenciais em incidentes críticos.
\end{itemize}

% Formação e Prêmios
\section{Formação \& Histórico}
\textbf{FGV --- Ciência de Dados \& Economia (2020--2022)} \hfill Rio de Janeiro, Brasil \\
Bolsista integral OBMEP; AED: 9,71. Saiu após dois anos para retornar ao Sul --- a bolsa exigia residência integral a 1.000 km da família; optou pelo trabalho prático em vez de concluir a graduação no Rio.

\vspace{3pt}
\textbf{IFPR --- Análise e Desenvolvimento de Sistemas (2023--2024)} \hfill União da Vitória, Brasil \\
Conceito A em Matemática para Computação. Saiu para se dedicar integralmente ao trabalho de engenharia de produção; a atuação em incidentes reais acelerava o aprendizado mais do que o currículo noturno.

\vspace{3pt}
\textbf{Medalha de Prata --- Olimpíada Brasileira de Matemática (OBMEP, 2016)} \hfill Nacional \\
Top 0,01\% entre 18M+ participantes. Medalha de Bronze (2014); 2x Menções Honrosas.

\end{document}
